% Source: https://artofproblemsolving.com/community/c3h2510871_is_the_implication_true
% Stable identity: aops-topic-2510871-q01
In mathematics, an implication ``if $P$, then $Q$'' is false only when $P$ is true and $Q$ is false; in every other case it is true.
\begin{enumerate}
\item[(a)] Decide whether each implication is true or false:
(i) If $2+2=3$, then $2+2=4$. \quad (ii) If $2+2=4$, then $2+2=3$. \quad (iii) If $2+2=3$, then $2+2=5$. \quad (iv) If $2+2=4$, then $2\cdot2=4$.
\item[(b)] Explain why this convention is the right one, using the statement ``for every integer $n$, if $n$ is divisible by $4$, then $n$ is even'' with $n=2$ and $n=3$.
\end{enumerate}
